% SHARD-ID: CA-07-EXCHANGE-STATISTICS-LAYER
% SHARD-TITLE: Exchange Symmetry and Statistics as Sector Selection
% SHARD-SUMMARY: Applies the sector machinery to the permutation action on tensor slots, recovering symmetric and antisymmetric Fock sectors from the full distinguishable carrier.
% SHARD-SUMMARY: Keeps internal symmetries separate from exchange symmetry and records how the two actions coexist on n-particle sectors.
% SHARD-SUMMARY: Defers para-statistics and genuinely anyonic braid statistics until their source, convention, and projector data are fixed.
% SHARD-KEYWORDS: exchange symmetry, statistics, symmetric Fock, antisymmetric Fock, permutation group, sector projector, anyonic deferral

\section{Exchange Symmetry and Statistics as Sector Selection}
\label{sec:exchange-statistics-layer}

\subsection{Sources and Dependencies}

This shard discharges part of the deferral in CA-03: ordinary Bose/Fermi
statistics are treated as sector selections on the full distinguishable Fock
carrier.  It depends on CA-05 for internal symmetry propagation and CA-06 for
projection/subspace/quotient semantics.

% Source: literature/md/1910.10619/1910.10619.md:39
%   Fock presentations often mix exchange symmetry and indefinite particle number.
% Source: literature/md/1910.10619/1910.10619.md:41
%   Bose-/Fermi-Fock space can be obtained by imposing exchange equivalence.
% Source: literature/md/2010.11121/2010.11121.md:258
%   bosonic Fock space uses symmetric tensor powers.
% Source: references/cft/Schottenloher2008/Schottenloher2008.md:1642
%   fermionic Fock/spinor space is built from exterior powers.
% Source: references/cft/Schottenloher2008/Schottenloher2008.md:6672
%   symmetric powers can also be presented as a quotient of tensor algebra.
% Source: references/text/TrebstShortIntroductionFibonacciAnyons.txt:42
%   two-dimensional anyon exchange can be Abelian phases or non-Abelian matrices.

\subsection{Two Symmetry Layers}

There are two distinct actions on many-particle carriers:
\[
  \text{internal symmetry }G,\qquad
  \text{exchange symmetry }S_n.
\]
If \(U:G\to U(H)\) is an internal one-particle representation, then on the
\(n\)-particle distinguishable sector \(H^{\hotimes n}\) it acts diagonally:
\[
  U_n(g)=U(g)^{\hotimes n}.
\]
The permutation group \(S_n\) acts by permuting tensor slots:
\[
  \Pi_n(\sigma)(v_1\otimes\cdots\otimes v_n)
  =
  v_{\sigma^{-1}(1)}\otimes\cdots\otimes v_{\sigma^{-1}(n)}.
\]
The two actions commute:
\[
  \Pi_n(\sigma)U_n(g)=U_n(g)\Pi_n(\sigma).
\]
Thus internal charge sectors and exchange-statistics sectors can be imposed in
either order in the ordinary tensor-slot setting.

\subsection{Bose and Fermi Projectors}

For finite \(S_n\), CA-06 gives averaging projectors.  The symmetric projector is
\[
  P_{\mathrm s,n}
  =
  \frac1{n!}\sum_{\sigma\in S_n}\Pi_n(\sigma),
\]
and the antisymmetric projector is
\[
  P_{\mathrm a,n}
  =
  \frac1{n!}\sum_{\sigma\in S_n}\operatorname{sgn}(\sigma)\Pi_n(\sigma).
\]
Their ranges are the symmetric and antisymmetric tensor powers:
\[
  \Ran(P_{\mathrm s,n})=H^{\hotimes_s n},\qquad
  \Ran(P_{\mathrm a,n})=\bigwedge^n H.
\]
The corresponding Fock sectors are
\[
  \Gamma_{\mathrm s}(H)=\hoplus_{n\ge0}\Ran(P_{\mathrm s,n}),\qquad
  \Gamma_{\mathrm a}(H)=\hoplus_{n\ge0}\Ran(P_{\mathrm a,n}).
\]
This makes precise the CA-03 statement that full/distinguishable Fock is the
carrier and Bose/Fermi spaces are later sector choices.

\subsection{Quotient Presentations}

The symmetric sector can also be presented as a quotient of the full tensor
algebra by the closed linear span of relations
\[
  v_1\otimes\cdots\otimes v_n
  -
  v_{\sigma^{-1}(1)}\otimes\cdots\otimes v_{\sigma^{-1}(n)}.
\]
The antisymmetric sector uses the signed relations
\[
  v_1\otimes\cdots\otimes v_n
  -
  \operatorname{sgn}(\sigma)
  v_{\sigma^{-1}(1)}\otimes\cdots\otimes v_{\sigma^{-1}(n)}.
\]
In the compiler, these quotient presentations are secondary: the primary sector
object is the orthogonal projector, because it fixes the inherited Hilbert
structure and the compressed observable algebra.

\subsection{Para and Anyonic Boundaries}

Para-statistics would require selecting other irreducible \(S_n\)-isotypic
components, hence central idempotents or matrix-block projectors for the group
algebra of \(S_n\).  This shard records that as a proposal, not as a completed
construction; a later shard must source the representation convention and the
projector normalization before using it.

Genuinely anyonic exchange is not an \(S_n\)-projection problem on ordinary tensor
slots.  It involves braid-group or categorical data and, in this project, must be
handled together with the fusion-category conventions for associators, braiding,
and \(F/R\)-symbols.  No braid-statistics projector is defined here.

\subsection{Lamport Dependency Skeleton}

{\small
\[
\setlength{\arraycolsep}{2pt}
\begin{array}{rcl}
D21 &:& U_n(g)=U(g)^{\hotimes n}\quad\text{(internal symmetry on }H^{\hotimes n}).\\
D22 &:& \Pi_n:S_n\to U(H^{\hotimes n})\quad\text{permutes tensor slots}.\\
D23 &:& P_{\mathrm s,n}=n!^{-1}\sum_{\sigma}\Pi_n(\sigma),\quad
        P_{\mathrm a,n}=n!^{-1}\sum_{\sigma}\operatorname{sgn}(\sigma)\Pi_n(\sigma).\\
L3 &:& \Pi_n(\sigma)U_n(g)=U_n(g)\Pi_n(\sigma).\\
L4 &:& P_{\mathrm s,n},P_{\mathrm a,n}\text{ are orthogonal projections by CA-06 finite averaging.}\\
T13 &:& \Gamma_{\mathrm s},\Gamma_{\mathrm a}\text{ are sector selections inside full Fock.}\\
T14 &:& \text{Anyonic exchange remains deferred to categorical braid data.}
\end{array}
\]
}

\subsection{Compiler Consequence}

The compiler now has its first nontrivial sector family:
\[
  \mathsf{Compile}(\operatorname{Bose}(\Fock(H)))
  =
  (\Gamma_{\mathrm s}(H),\ P_{\mathrm s},\ \text{chosen compressed algebra},\ldots)
\]
and similarly for \(\operatorname{Fermi}\).  The sector constructor is separate
from Fock itself, preserving the separation of indefinite particle number from
exchange symmetry.
